% !PRINTED 372
\weberpage{372}

Die Substitution (5) heisst aus (1) und (4) \gesperrt{zusammengesetzt}.

Es macht sich das Bedürfniss nach einer abgekürzten Bezeichnung
 dieser Substitutionen geltend. Da es häufig nicht auf
die Variablen, sondern nur auf die Substitutionszahlen $\alpha, \beta, \gamma, \delta$
ankommt, so bezeichnet man die ganze Substitution (1) durch
die Substitutionszahlen oder auch nur durch einen einfachen
Buchstaben
\weq{9}{S = \begin{pmatrix} \alpha, & \beta \\ \gamma, & \delta \end{pmatrix},}

und schreibt dann, wenn es nöthig ist, die Gleichung (1) so:
\[
y = S(x) = \begin{pmatrix} \alpha, & \beta \\ \gamma, & \delta \end{pmatrix}(x).
\]

Die Zusammensetzung zweier Substitutionen bezeichnet man
durch Nebeneinandersetzen der Zeichen, wobei aber auf die
\gesperrt{Reihenfolge} zu achten ist, also
\weq{10}{\begin{pmatrix} \alpha, & \beta \\ \gamma, & \delta \end{pmatrix} \begin{pmatrix} \alpha', & \beta' \\ \gamma', & \delta' \end{pmatrix} = \begin{pmatrix} \alpha'', & \beta'' \\ \gamma'', & \delta'' \end{pmatrix},}

oder
\weq{11}{S S' = S''.}

Die Formeln (6) enthalten die Vorschrift, nach der eine
zusammengesetzte Substitution zu bilden ist. Man kann sie aus
der Multiplicationsregel der Determinanten ableiten, muss aber
beachten, dass im ersten Factor nach Zeilen, im zweiten nach
Colonnen summirt werden muss, wie eben die Formeln (6) zeigen.

Ebenso wie man zwei Substitutionen zusammensetzt, kann
man auch die Zusammensetzung von mehreren bilden:
\[
S S_1 S_2 S_3 \ldots
\]

Bei der Bildung der zusammengesetzten Substitution dürfen,
im Allgemeinen wenigstens, die Componenten nicht vertauscht
werden; wohl aber kann man nach Belieben zwei benachbarte
zu einer zusammenfassen, dann wieder zwei u. s. f., d. h. es gilt
zwar nicht das \gesperrt{commutative}, wohl aber das \gesperrt{associative}
Gesetz.

Dies folgt unmittelbar daraus, dass, wenn $x$ durch $x_1$, $x_1$
durch $x_2$, $x_2$ durch $x_3$ ausgedrückt ist, der Ausdruck von $x$ durch
$x_3$ entweder dadurch gefunden werden kann, dass man zuerst
$x_1$ in $x$ durch $x_2$ und dann $x_2$ durch $x_3$ ausdrückt, oder dass
%CONT>>
